\section{Selection, omitted evidence and diagnostic limits}

\begin{annexitem}{annex:p15.missing.selection}{thm:p15-missing-selection}{Conditional availability, recovery assumptions and nonidentification}
\begin{proof}
For each support value \(y\), the joint probability of that result
and availability is \(p(y)s(y)\). Summing these probabilities gives
\(\Pr(R=1)=\rho>0\). Conditional division therefore yields
\(q(y)=p(y)s(y)/\rho\). Its entries are nonnegative and sum to
one. Multiplying by \(y\) and summing gives
\[
 E(Y\mid R=1)-E(Y)
 =\frac{E[Ys(Y)]-E(Y)E[s(Y)]}{\rho},
\]
which is the stated covariance formula. This does not assume a common
clinical effect or that every numerical estimate is unbiased for it;
it describes the distribution of results entering the synthesis.

If every specified \(s(y)>0\), the conditional-law identity implies
\(q(y)/s(y)=p(y)/\rho\). The sum over \(y\) is \(1/\rho\),
so normalizing recovers the full distribution. Conditional on a result
\(y\), the expected contribution \(RY/s(Y)\) is
\(s(y)y/s(y)=y\). Summing under \(p\) gives its expectation
\(E(Y)\). This expectation is over the full instance and selection
law. Averaging inverse-weighted values only over reported results without
the corresponding normalization is a different calculation; a realized
ratio estimator is not thereby proved unbiased at finite sample size.

In the constructed two-value example,
\(\rho=(1/2)(1/3)+(1/2)1=2/3\).
Thus \(q(-1)=1/4\), \(q(1)=3/4\), and its mean is
\(-1/4+3/4=1/2\), whereas the full symmetric mean is zero.
The no-selection model with these same \(q\) probabilities has
\(s=1\), \(\rho=1\), and the same observed law. Its full
mean is \(1/2\). Hence neither the full mean nor the selection
mechanism is identified by the available distribution alone.

For a zero-selection-probability boundary, suppose the only reported
result is one. For any \(0<\eta\leq1\), give that full result
mass \(\eta\) with selection probability one, and give an unreported
result \(a\ne1\) mass \(1-\eta\) with selection probability zero.
Conditional on availability the result is always one, independently
of \(a,\eta\). The full mean \(\eta+(1-\eta)a\) can change
without any change to the observed law. Division by its zero selection
probability is undefined and cannot recover that mass. This is the
support failure that a positive-selection assumption excludes.
\end{proof}
\end{annexitem}

\begin{annexitem}{annex:p15.missing.cohort}{thm:p15-missing-cohort}{Closed-frame accounting and the excluded target weight}
\begin{proof}
A complete enumeration of the declared inception cohort lists every
study in \(G\), so its cohort-only eligible studies are accounted for.
This premise concerns identification; it makes no statement that the
record of any one listed study contains its eligible effect or precision.
Nor is it established merely by searching a register. Completeness of
the frame and its enumeration are the explicit premises.

For \(0<\lambda<1\), define
\[
 \theta_G=\frac{\sum_{i\in G}b_i\theta_i}{\lambda},\qquad
 \theta_{G^c}=\frac{\sum_{i\notin G}b_i\theta_i}{1-\lambda}.
\]
Splitting the full weighted sum into these two disjoint parts gives
the stated decomposition. The normalized weights on the complement
are positive and sum to one, so \(L\leq\theta_{G^c}\leq U\).
Multiplying by the positive \(1-\lambda\) gives the two full-target
bounds. Set every complement effect equal to \(L\) to attain the
lower endpoint, or equal to \(U\) to attain the upper endpoint.
Both choices leave the entire cohort and its observations unchanged.
If \(\lambda=1\), the complement has zero weight and the full
weighted target is \(\theta_G\).

With a proper cohort and unspecified complementary effects, take two
different common values on its complement. The cohort information is
unchanged, while the full target changes by \((1-\lambda)\) times
their difference. This proves the nonidentification claim. If the
full target weights or their complementary mass are themselves unknown,
the numerical interval cannot be asserted without those further inputs.
All statements concern the defined effects and weights. Sampling
uncertainty in estimating \(\theta_G\) would require an additional
confidence or posterior calculation; it is not supplied by these bounds.
\end{proof}
\end{annexitem}

\begin{annexitem}{annex:p15.missing.bounds}{thm:p15-missing-weight-bounds}{Every range endpoint and the exact threshold calculation}
\begin{proof}
The available weighted numerator is \(W_oM_o\). Adding the
missing results gives
\[
        M=\frac{W_oM_o+\sum_jw_jy_j}{W_o+W_m}.
\]
Its denominator is fixed and strictly positive because \(W_o>0\).
Each missing coordinate has coefficient \(w_j/(W_o+W_m)>0\).
It is therefore nondecreasing in each \(y_j\); setting all
coordinates to their lower bounds attains the minimum and setting
all to their upper bounds attains the maximum. Intermediate possible
values cannot leave this interval. If there are no missing coordinates,
the interval is the single point \(M_o\).

When \(W_m>0\), substitute
\(\sum_jw_jy_j=W_mM_m\) and subtract \(M_o\). The
result is \(W_m(M_m-M_o)/(W_o+W_m)\). To compare \(M\)
with threshold \(d\), multiply by its positive denominator
and move \(W_oM_o\) to the other side. This gives exactly
the displayed inequality, without changing its direction.

In the example, \(W_o=3\), \(M_o=2\), \(W_m=2\).
The two omitted values can have sum anywhere in \([-8,0]\).
The full mean is therefore \((6+y_1+y_2)/5\), with range
\([-2/5,6/5]\). A nonpositive mean is equivalent to
\(y_1+y_2\leq-6\), including equality at zero.
The originally available mean two is compatible with each of
these completed datasets. Selecting a summary formula for the
available data does not reveal which completion occurred.

The calculation extends fixed-weight accounting in
\cref{thm:p15-synthesis-sensitivity}; it does not assume missing
standard errors are known from participant counts. Re-estimating
heterogeneity, variances or outcome conversions changes its weights
and requires recomputation under each sensitivity scenario.
It bounds deterministic completed numerical summaries, rather than
claiming a probability of coverage for an unobserved population effect.
\end{proof}
\end{annexitem}

\begin{annexitem}{annex:p15.missing.small-weight}{thm:p15-missing-small-weight}{Relative weighting and what reweighting cannot recover}
\begin{proof}
The two weights under allowance \(u\) are
\(1/(v_L+u)\) and \(1/(v_S+u)\), as in
\cref{def:p15-synthesis-random-model}. Multiplying the weighted
ratio's numerator and denominator by
\((v_L+u)(v_S+u)\) gives the expression for \(M(u)\).
Differentiating that ratio gives numerator
\[
 (y_L+y_S)(v_L+v_S+2u)
 -2[(v_S+u)y_L+(v_L+u)y_S]
       =(v_S-v_L)(y_S-y_L).
\]
The denominator is \((v_L+v_S+2u)^2>0\). Dividing numerator
and denominator of \(M\) by \(u\) proves its limit as
\(u\to\infty\). The normalized coefficient of \(y_S\)
is \((v_L+u)/(v_L+v_S+2u)\), with derivative
\((v_S-v_L)/(v_L+v_S+2u)^2>0\). Thus the less precise
result gains relative weight; the sign of the overall mean shift
also depends on which estimate is larger.

The stated numerical inputs give
\(M(u)=3(1+u)/(10+2u)\). At zero this is \(3/10\),
and at four it is \(15/18=5/6\). Finally any weighted mean
of finitely many strictly positive available values with positive
weights is strictly positive: its numerator and denominator are
both positive. This remains true at every finite allowance and
at the equal-weight limit. Hence such reweighting cannot produce
the negative completed mean allowed by the preceding missing-result
example. Modifying the variance model addresses relative precision,
not recovery of an unavailable numerical value.
\end{proof}
\end{annexitem}

\begin{annexitem}{annex:p15.missing.funnel}{thm:p15-missing-funnel}{Positive-probability observations under three distinct causes}
\begin{proof}
Let each displayed estimate stand for an interval of half-width
\(\epsilon\), with \(0<\epsilon<1/40\). This keeps the
five precise displayed intervals distinct, and places every
less-precise interval strictly above one. At any real point, a
normal density with positive variance and any finite mean is
strictly positive. Its integral over an interval of positive
length is therefore strictly positive.

Under model 1, the probability of all ten interval events is the
product of ten positive normal probabilities, by independence.
It is positive even though the less-precise positive estimates
are unusual under a zero true effect. Every eligible study is
available, so selective reporting is absent in this model.

Under model 2, replace only the five less-precise means by two.
Their interval probabilities remain positive, as do those of
the precise zero-mean studies. Their independent product is
again positive. Every eligible study is still available, but
the study types now have different true effects. This is a
heterogeneity model, not a missing-study model.

For model 3, a zero-mean, unit-SE candidate result is available
only when it exceeds one. Conditional on availability its density
is the positive normal density divided by \(\Pr(Y>1)>0\)
on that half-line, and zero elsewhere. Each of the five displayed
less-precise intervals lies inside that half-line, so each has
positive conditional probability. The five precise studies are
unselected zero-mean normal results and also have positive interval
probabilities. Independent retained results thus have positive
probability of the displayed plot.

This can also be realized with a finite candidate set: generate
seven independent less-precise zero-mean studies, have five fall
in the five displayed intervals and the other two in \((-1,0)\).
That event has positive probability. The rule makes the former
five available and withholds the latter two, while all five
precise studies remain available. The observed ten studies then
give the same rounded plot as in the two no-selection models.

Positive probability establishes compatibility; it does not
assert equal probabilities, equal likelihoods, equal prior
plausibility or an inability to gather evidence discriminating
these mechanisms. In particular external registry or protocol
information can rule out the withheld studies in this specific
construction. The conclusion is that the plotted asymmetry by
itself does not logically determine a unique publication cause.
\end{proof}
\end{annexitem}
